\documentclass[11pt,reqno]{amsart}

\usepackage[T1]{fontenc}
\usepackage{newtxtext,newtxmath}
\let\Bbbk\relax
\usepackage{amsmath,amssymb,amsfonts}
\usepackage{microtype}
\usepackage{xcolor}
\usepackage[numbers]{natbib}
\usepackage[colorlinks=true,linkcolor=blue,citecolor=blue!60!black,urlcolor=blue!60!black]{hyperref}

\newenvironment{contextnote}
{\par\medskip\noindent\hrulefill\par\smallskip\small}
{\par\smallskip\noindent\hrulefill\par\medskip}

\begin{document}

\title{History-dependent adaptive response geometry}
\maketitle

\section{Results}

The law of requisite variety relates successful regulation to the range of responses available to a regulator \cite{ashby1956}. Related ideas appear in ecological response diversity \cite{elmqvist2003} and in multivariate evolutionary response, where a covariance matrix determines the directions along which selection can move a population \cite{lande1979}. Here we use a minimal adaptive model to make a different part of this picture explicit: adaptation changes the distribution of available responses, and therefore changes the geometry governing subsequent adaptation.

\subsection{Adaptive response model}

Consider a finite repertoire of response vectors
\begin{equation}
\mathbf a_i\in\mathbb R^d,
\qquad
i=1,\ldots,n,
\end{equation}
with nonnegative weights
\begin{equation}
p_i\geq0,
\qquad
\sum_i p_i=1.
\label{eq:simplex}
\end{equation}
The realized response is
\begin{equation}
\mathbf r
=
\sum_i p_i\mathbf a_i.
\label{eq:response}
\end{equation}
The vectors $\mathbf a_i$ are fixed. Adaptation only redistributes weight among responses that are already available.

Let $\mathbf y(t)\in\mathbb R^d$ be the response required by the current environment and define
\begin{equation}
E
=
\frac12\|\mathbf r-\mathbf y\|^2.
\label{eq:error}
\end{equation}
We let the weights descend this mismatch on the simplex using the Shahshahani, or replicator, gradient flow \cite{shahshahani1979},
\begin{equation}
\dot p_i
=
-p_i
\left(
\frac{\partial E}{\partial p_i}
-
\sum_jp_j\frac{\partial E}{\partial p_j}
\right).
\end{equation}
Since
\begin{equation}
\frac{\partial E}{\partial p_i}
=
\mathbf a_i^{\mathsf T}(\mathbf r-\mathbf y),
\end{equation}
the adaptive law is
\begin{equation}
\dot p_i
=
p_i
(\mathbf a_i-\mathbf r)^{\mathsf T}
(\mathbf y-\mathbf r).
\label{eq:replicator}
\end{equation}
The overall adaptation rate has been absorbed into time. A response whose direction agrees with the current mismatch gains weight relative to the mean response.

Define the weighted covariance
\begin{equation}
\mathbf C
=
\sum_i p_i
(\mathbf a_i-\mathbf r)
(\mathbf a_i-\mathbf r)^{\mathsf T}.
\label{eq:covariance}
\end{equation}
Differentiating Eq.~\eqref{eq:response} gives the exact relation
\begin{equation}
\boxed{
\dot{\mathbf r}
=
\mathbf C(\mathbf y-\mathbf r)
}.
\label{eq:response-dynamics}
\end{equation}
Thus the covariance of the currently represented responses is also the operator that converts environmental mismatch into adaptive motion.

For a fixed environment,
\begin{equation}
\dot E
=
-(\mathbf r-\mathbf y)^{\mathsf T}
\mathbf C
(\mathbf r-\mathbf y)
\leq0.
\label{eq:error-decay}
\end{equation}
At a matched state, a small change $\delta\mathbf y$ gives to first order
\begin{equation}
\delta\dot{\mathbf r}
=
\mathbf C
(\delta\mathbf y-\delta\mathbf r).
\label{eq:linear-response}
\end{equation}
Hence the eigenvectors of $\mathbf C$ are the local directions of adaptation and its nonzero eigenvalues set the corresponding relaxation rates. The role of covariance as a response operator is classical in several settings; the point here is that $\mathbf C$ is itself changed by the same redistribution process that produces adaptation.

\subsection{An exact history coordinate}

The $n$ weights can be written in terms of a single $d$-dimensional history variable. Define
\begin{equation}
\boldsymbol\theta(t)
=
\int_0^t
\left[
\mathbf y(s)-\mathbf r(s)
\right]
\,\mathrm ds.
\label{eq:theta}
\end{equation}
Integrating Eq.~\eqref{eq:replicator} gives
\begin{equation}
\boxed{
p_i(t)
=
\frac{
p_i(0)
e^{\mathbf a_i^{\mathsf T}\boldsymbol\theta(t)}
}{
\sum_j
p_j(0)
e^{\mathbf a_j^{\mathsf T}\boldsymbol\theta(t)}
}
}.
\label{eq:weights-solution}
\end{equation}
Introduce the log-partition function
\begin{equation}
\psi(\boldsymbol\theta)
=
\log
\sum_i
p_i(0)
e^{\mathbf a_i^{\mathsf T}\boldsymbol\theta}.
\label{eq:psi}
\end{equation}
Standard exponential-family identities then give \cite{amari2016}
\begin{equation}
\boxed{
\mathbf r
=
\nabla\psi,
\qquad
\mathbf C
=
\nabla^2\psi
}.
\label{eq:gradient-hessian}
\end{equation}
The full dynamics therefore reduces exactly to
\begin{equation}
\boxed{
\dot{\boldsymbol\theta}
=
\mathbf y(t)
-
\nabla\psi(\boldsymbol\theta)
}.
\label{eq:theta-dynamics}
\end{equation}
The environmental history does not enter through time itself, but through the accumulated mismatch $\boldsymbol\theta$. Once $\boldsymbol\theta$ is known, both the current response and the local response operator follow from the same convex function $\psi$.

For a constant environment $\mathbf y$, Eq.~\eqref{eq:theta-dynamics} is the gradient flow
\begin{equation}
\dot{\boldsymbol\theta}
=
-\nabla V_{\mathbf y},
\qquad
V_{\mathbf y}(\boldsymbol\theta)
=
\psi(\boldsymbol\theta)
-
\mathbf y^{\mathsf T}\boldsymbol\theta.
\label{eq:dual-potential}
\end{equation}
This makes the long-time behavior a geometric problem.

\subsection{Reachable response geometry}

Because the weights remain on the simplex, the set of responses obtainable by redistribution is
\begin{equation}
\mathcal R
=
\operatorname{conv}
\{\mathbf a_1,\ldots,\mathbf a_n\}.
\label{eq:repertoire}
\end{equation}
The position of $\mathbf y$ relative to $\mathcal R$ separates three qualitatively different regimes.

\begin{contextnote}
\textit{Geometric picture.}
The convex hull $\operatorname{conv}\{\mathbf a_i\}$ is the set obtained by taking every weighted average of the response vectors with nonnegative weights that sum to one. In two dimensions it is the polygon whose corners are the outermost $\mathbf a_i$. Its relative interior means the interior inside the smallest affine space containing that polygon. A face is a boundary piece touched by a supporting line or hyperplane: an edge or a vertex in two dimensions, and their higher-dimensional analogues.
\end{contextnote}

If
\begin{equation}
\mathbf y\in\operatorname{relint}\mathcal R,
\end{equation}
there is a finite $\boldsymbol\theta_*$ satisfying
\begin{equation}
\nabla\psi(\boldsymbol\theta_*)
=
\mathbf y.
\end{equation}
For positive initial weights, every response retains positive asymptotic weight and
\begin{equation}
\mathbf C(t)
\longrightarrow
\mathbf C_*
=
\nabla^2\psi(\boldsymbol\theta_*).
\label{eq:interior-limit}
\end{equation}
On the affine span of the repertoire, $\mathbf C_*$ is nonsingular. Adaptation reaches the target without asymptotically suppressing a direction of the represented repertoire.

The boundary is different. Let $\mathbf y$ lie in the relative interior of a proper face $F\subset\mathcal R$. Any finite $\boldsymbol\theta$ in Eq.~\eqref{eq:weights-solution} gives positive weight to every initially represented response and therefore places $\mathbf r$ in the relative interior of $\mathcal R$. Reaching a boundary target consequently requires
\begin{equation}
\|\boldsymbol\theta(t)\|
\longrightarrow
\infty,
\end{equation}
while the weights of responses outside $F$ vanish.

For a codimension-one face, choose a unit normal $\mathbf n$ and $h$ such that
\begin{equation}
\mathbf n^{\mathsf T}\mathbf a_i
=
h
\quad
\text{for } \mathbf a_i\in F,
\end{equation}
and define the positive gaps
\begin{equation}
g_i
=
h-\mathbf n^{\mathsf T}\mathbf a_i
>0
\quad
\text{for } \mathbf a_i\notin F.
\label{eq:gaps}
\end{equation}
Let
\begin{equation}
g
=
\min_{\mathbf a_i\notin F}g_i.
\end{equation}
Writing
\begin{equation}
s(t)
=
\mathbf n^{\mathsf T}\boldsymbol\theta(t),
\end{equation}
the normal component of Eq.~\eqref{eq:theta-dynamics} is
\begin{equation}
\dot s
=
\mathbf n^{\mathsf T}(\mathbf y-\mathbf r)
=
\sum_{\mathbf a_i\notin F}
p_i g_i.
\label{eq:s-normal}
\end{equation}
Once the tangential coordinates have relaxed inside the face, Eq.~\eqref{eq:weights-solution} gives
\begin{equation}
p_i
\asymp
e^{-g_i s}
\quad
\text{outside }F.
\label{eq:boundary-weight-s}
\end{equation}
The smallest gap dominates Eq.~\eqref{eq:s-normal}, so
\begin{equation}
\dot s
\asymp
e^{-gs},
\qquad
s(t)
=
\frac{1}{g}\log t+O(1).
\label{eq:boundary-s}
\end{equation}
Consequently,
\begin{equation}
\boxed{
p_i(t)
\asymp
t^{-g_i/g}
}
\label{eq:boundary-weights}
\end{equation}
for responses outside the selected face. In particular, the least-separated responses decay as $t^{-1}$.

The same law appears directly in the response operator. Along the normal direction,
\begin{equation}
\mathbf n^{\mathsf T}\mathbf C\mathbf n
=
\operatorname{Var}_{p}
\left(
\mathbf n^{\mathsf T}\mathbf a
\right),
\end{equation}
and therefore
\begin{equation}
\boxed{
\mathbf n^{\mathsf T}\mathbf C(t)\mathbf n
\asymp
\frac{1}{t}
}.
\label{eq:boundary-capacity}
\end{equation}
Adaptation to a boundary target is exact, but it produces an algebraic loss of response capacity transverse to the face on which the target lies.

\begin{contextnote}
\textit{Why a face matters.}
A face collects responses that are equally extreme in some direction $\mathbf n$. If the environment demands a point on that face, only mixtures of responses on the face are needed asymptotically. Responses lying behind it acquire smaller and smaller weights. The gap $g_i$ is simply their distance from the supporting hyperplane measured along $\mathbf n$.
\end{contextnote}

Finally, consider
\begin{equation}
\mathbf y\notin\mathcal R.
\end{equation}
The best achievable response is the Euclidean projection
\begin{equation}
\mathbf r_*
=
\Pi_{\mathcal R}\mathbf y,
\label{eq:projection}
\end{equation}
and the residual
\begin{equation}
\mathbf e
=
\mathbf y-\mathbf r_*
\end{equation}
is normal to a supporting face of $\mathcal R$ at $\mathbf r_*$ \cite{rockafellar1970}. Since $\mathbf r(t)\to\mathbf r_*$,
\begin{equation}
\boldsymbol\theta(t)
=
\mathbf e\,t+o(t).
\label{eq:outside-theta}
\end{equation}
Let
\begin{equation}
F_*
=
\operatorname*{arg\,max}_{\mathbf a_i}
\mathbf e^{\mathsf T}\mathbf a_i
\end{equation}
be the exposed face selected by the residual mismatch. For a response outside this face define
\begin{equation}
\Delta_i
=
\mathbf e^{\mathsf T}
(\mathbf r_*-\mathbf a_i)
>0.
\label{eq:external-gap}
\end{equation}
Equation~\eqref{eq:weights-solution} then gives
\begin{equation}
\boxed{
p_i(t)
=
e^{-\Delta_i t+o(t)}
}.
\label{eq:outside-weights}
\end{equation}
An unattainable environment therefore suppresses unused response directions exponentially rather than algebraically.

The three regimes can be summarized as
\begin{equation}
\begin{array}{ccl}
\mathbf y\in\operatorname{relint}\mathcal R
&\Longrightarrow&
\text{finite limiting response operator},
\\[3pt]
\mathbf y\in\partial\mathcal R
&\Longrightarrow&
\text{algebraic specialization},
\\[3pt]
\mathbf y\notin\mathcal R
&\Longrightarrow&
\text{exponential specialization}.
\end{array}
\label{eq:three-regimes}
\end{equation}

\subsection{Past specialization sets future recovery costs}

The suppression laws above become measurable when the environment changes. Suppose the system is exposed to one environment until time $T$, after which the target changes to $\mathbf y^+$. Consider a response $i$ whose weight has become small. Under the new environment,
\begin{equation}
\frac{\mathrm d}{\mathrm dt}
\log p_i
=
(\mathbf a_i-\mathbf r)^{\mathsf T}
(\mathbf y^+-\mathbf r).
\label{eq:post-switch-growth}
\end{equation}
Assume that while $p_i$ grows from its value $p_i(T)$ to a fixed small threshold $\eta$, its growth advantage remains bounded,
\begin{equation}
0<\sigma_-
\leq
(\mathbf a_i-\mathbf r)^{\mathsf T}
(\mathbf y^+-\mathbf r)
\leq
\sigma_+.
\label{eq:growth-bounds}
\end{equation}
If $\tau_i(T)$ is the time required to reach $p_i=\eta$, integration of Eq.~\eqref{eq:post-switch-growth} gives
\begin{equation}
\boxed{
\frac{1}{\sigma_+}
\log\frac{\eta}{p_i(T)}
\leq
\tau_i(T)
\leq
\frac{1}{\sigma_-}
\log\frac{\eta}{p_i(T)}
}.
\label{eq:recovery-bound}
\end{equation}
The cost of recovering a previously suppressed response is therefore set by how strongly the previous environment depleted it.

Combining Eq.~\eqref{eq:recovery-bound} with the long-time specialization laws gives
\begin{equation}
\boxed{
\begin{array}{ccl}
p_i(T)\to p_i^*>0
&\Longrightarrow&
\tau_i(T)=O(1),
\\[3pt]
p_i(T)\asymp T^{-\alpha_i}
&\Longrightarrow&
\tau_i(T)=\Theta(\log T),
\\[3pt]
p_i(T)=e^{-\Delta_iT+o(T)}
&\Longrightarrow&
\tau_i(T)=\Theta(T).
\end{array}
}
\label{eq:recovery-scaling}
\end{equation}
Thus two environments that are both followed for a long time can leave qualitatively different memories in the same repertoire. Adaptation to an attainable boundary target creates a logarithmically growing recovery delay, whereas persistent forcing beyond the available repertoire creates a recovery delay proportional to the duration of the previous exposure.

\begin{contextnote}
\textit{Interpretation of the scaling.}
Nothing has been deleted from the repertoire: every $p_i$ remains positive at finite time. The memory is instead stored in relative abundance. Recovering a rare response requires multiplying its weight back to an appreciable level. Polynomial suppression therefore costs only a logarithm of the previous exposure time, while exponential suppression turns the previous exposure time itself into a future recovery timescale.
\end{contextnote}

Equations~\eqref{eq:weights-solution}, \eqref{eq:boundary-weights}, \eqref{eq:outside-weights}, and \eqref{eq:recovery-scaling} give a direct chain from environmental history to future adaptation:
\begin{equation}
\boxed{
\text{environmental history}
\longrightarrow
\boldsymbol\theta
\longrightarrow
\{p_i\},\,\mathbf C
\longrightarrow
\text{future recovery time}.
}
\end{equation}
The fixed repertoire is essential to this result. The model describes selection and retention among existing responses, not the construction of new response vectors. It therefore provides a baseline for a later extension in which the repertoire itself can change.

\bibliographystyle{unsrtnat}
\bibliography{refs}

\end{document}
